\documentclass[12pt,twoside,a4paper,openany]{article}

    \input{preambula_pakiety.tex}
    \input{preambula_ustawienia.tex}

    % polecenia zdefiniowane w pakiecie strona_tytulowa.sty
    \title{Wykrywanie obrazów generowanych przez sztuczną inteligencję}
    \author{Agnieszka Ryś}
    \authorI{}
    \authorII{}

    \uczelnia{AKADEMIA NAUK STOSOWANYCH \\W NOWYM SĄCZU}
    \instytut{Wydział Nauk Inżynieryjnych}
    \kierunek{Katedra Informatyki}
    \praca{DOKUMENTACJA PROJEKTOWA}
    \przedmiot{ANALIZA OBRAZÓW Z ELEMENTAMI SZTUCZNEJ INTELIGENCJI W DETEKCJI ZAGROŻEŃ}
    \prowadzacy{mgr inż. Daniel Drozd}
    \rok{2026}

\usepackage{fancyhdr}

\begin{document}
    \setlength{\headheight}{23pt}

    \renewcommand{\figurename}{Rys.}
    \renewcommand{\tablename}{Tab.}
    \thispagestyle{empty}
    \stronatytulowa

    \pagestyle{fancy}

    \newpage

    % formatowanie spisu treści i nagłówków
    \renewcommand{\cftbeforesecskip}{8pt}
    \renewcommand{\cftsecafterpnum}{\vskip 8pt}
    \renewcommand{\cftparskip}{3pt}
    \renewcommand{\cfttoctitlefont}{\Large\bfseries\sffamily}
    \renewcommand{\cftsecfont}{\bfseries\sffamily}
    \renewcommand{\cftsubsecfont}{\sffamily}
    \renewcommand{\cftsubsubsecfont}{\sffamily}
    \renewcommand{\cftparafont}{\sffamily}

    \tableofcontents
    \thispagestyle{fancy}
    \newpage

%%%%%%%%%%%%%%%%%%% treść główna dokumentu %%%%%%%%%%%%%%%%%%%%%%%%%

   \newpage
\section{Cel projektu}

Celem projektu jest stworzenie systemu wykrywania obrazów generowanych przez sztuczną inteligencję. System wykorzystuje metody uczenia maszynowego do analizy cech wizualnych i identyfikacji obrazów, które zostały wygenerowane przez modele generatywne. W kontekście współczesnych zagrożeń problem ten ma znaczenie praktyczne między innymi dla wykrywania dezinformacji, podszywania się pod inne osoby oraz obniżania wiarygodności obrazu jako nośnika informacji. W projekcie uwzględniono zarówno klasyczny etap \texttt{real vs fake}, jak i osobny etap twarzowy, przeznaczony do analizy portretów i potencjalnych deepfake'ów.

\subsection{Zakres projektu}

Zakres projektu obejmuje zaprojektowanie, implementację oraz weryfikację systemu, który klasyfikuje obrazy na prawdziwe i wygenerowane przez AI. Prace objęły przygotowanie danych, budowę modeli, analizę odporności oraz sprawdzenie zachowania modeli na danych spoza głównego rozkładu treningowego.

\begin{itemize}
    \item przygotowanie zbiorów danych i ich podział na zbiory treningowy, walidacyjny i testowy,
    \item implementacja modelu globalnego opartego o klasyfikację całego obrazu,
    \item implementacja osobnego etapu twarzowego z wykrywaniem twarzy i klasyfikacją cropów,
    \item przygotowanie benchmarku OOD dla nowszych generatorów obrazów,
    \item analiza skuteczności, odporności oraz ograniczeń modeli,
    \item przygotowanie lokalnego interfejsu demonstracyjnego do ręcznych testów i analizy wyników.
\end{itemize}

\subsection{Wymagania funkcjonalne}

\begin{itemize}
    \item wczytanie danych z ustalonej struktury katalogów \texttt{train/val/test},
    \item konfiguracja parametrów treningu i modelu z pliku konfiguracyjnego,
    \item preprocessing obrazów oraz wdrożenie augmentacji uodparniających system na popularne degradacje jakości,
    \item trening modelu oraz zapis najlepszego checkpointu,
    \item ewaluacja modelu i raport metryk \texttt{Accuracy}, \texttt{Precision}, \texttt{Recall}, \texttt{F1} i \texttt{ROC-AUC},
    \item automatyczne wydzielanie rejonu twarzy i podanie go do dedykowanego modelu twarzowego,
    \item wizualizacja przesłanek decyzyjnych modelu przy użyciu map Grad-CAM \cite{selvaraju2017gradcam}.
\end{itemize}

\subsection{Wymagania niefunkcjonalne}

\begin{itemize}
    \item przenośność --- możliwość uruchomienia na środowisku Windows/Linux,
    \item wydajność --- wsparcie dla akceleracji GPU i rozsądny czas trenowania,
    \item odtwarzalność --- zapisywanie konfiguracji, checkpointów i podsumowań eksperymentów,
    \item niezawodność --- czytelne zarządzanie wyjątkami, np. w przypadku braku wykrytej twarzy,
    \item interpretowalność --- możliwość ręcznej analizy logiki decyzyjnej modelu,
    \item bezpieczeństwo --- lokalne przetwarzanie danych bez konieczności wysyłania obrazów do usług zewnętrznych.
\end{itemize}

   \newpage
\section{Analiza problemu}

Rozwój modeli generatywnych spowodował, że obrazy syntetyczne stały się łatwe do wytworzenia i coraz trudniejsze do odróżnienia od zwykłych fotografii. Problem ten ma znaczenie praktyczne w kilku obszarach: dezinformacji, podszywaniu się pod inne osoby, publikowaniu sfałszowanych materiałów wizualnych oraz obniżaniu wiarygodności obrazu jako nośnika informacji. W projekcie przyjęto, że pierwszym krokiem powinno być zbudowanie modelu bazowego wykrywającego obrazy typu \texttt{real} oraz \texttt{fake}.

\subsection{Charakter problemu}

Obrazy generowane przez AI mogą zawierać artefakty trudne do uchwycenia dla człowieka, ale rozpoznawalne przez model uczony na odpowiednio dobranym zbiorze danych. Do takich cech można zaliczyć:

\begin{itemize}
    \item nienaturalną spójność lub niespójność tekstur,
    \item charakterystyczne błędy w drobnych detalach,
    \item zbyt gładkie przejścia tonalne,
    \item artefakty wynikające z procesu generacji lub późniejszej kompresji obrazu.
\end{itemize}

W praktyce problem nie sprowadza się jednak tylko do znalezienia jednego modelu o wysokiej skuteczności. Istotne jest również to, czy model pozostaje odporny na pogorszenie jakości obrazu, takie jak rozmycie, kompresja JPEG czy zmniejszenie rozdzielczości. Z punktu widzenia zastosowań praktycznych właśnie te przypadki są szczególnie ważne, ponieważ obrazy rozpowszechniane w mediach społecznościowych lub komunikatorach rzadko zachowują oryginalną jakość.

\subsection{Cel etapu bazowego}

W pierwszym etapie projektu założono budowę rozwiązania, które:

\begin{itemize}
    \item przyjmuje cały obraz jako wejście,
    \item klasyfikuje go do jednej z dwóch klas: \texttt{real} lub \texttt{fake},
    \item zapisuje pełne metryki i checkpointy,
    \item pozwala na wznowienie treningu,
    \item umożliwia interpretację decyzji modelu za pomocą Grad-CAM,
    \item może zostać później rozszerzone o osobny moduł analizy twarzy.
\end{itemize}

Takie podejście pozwala najpierw zbudować solidny punkt odniesienia, a dopiero później przejść do trudniejszego zagadnienia deepfake twarzy.

\subsection{Najważniejsze trudności}

W trakcie analizy problemu zidentyfikowano kilka kluczowych trudności:

\begin{itemize}
    \item obrazy spoza rozkładu treningowego mogą być klasyfikowane z dużą pewnością, ale błędnie,
    \item portrety niskiej jakości i obrazy po kompresji mogą powodować fałszywe alarmy,
    \item ilustracje i anime nie należą do głównego zakresu modelu fotograficznego,
    \item skuteczność mierzona na zbiorze testowym nie zawsze przekłada się wprost na zachowanie modelu na obrazach z zewnątrz.
\end{itemize}

Z tego względu poza samym treningiem konieczne było dodanie testu odporności oraz prostego mechanizmu ostrzegania o niskiej jakości obrazu w interfejsie demonstracyjnym.

   \newpage
\section{Projektowanie}

\subsection{Architektura rozwiązania}

Projekt został podzielony na kilka logicznych modułów, z których każdy odpowiada za inny fragment eksperymentu:

\begin{itemize}
    \item \texttt{dataset.py} --- przygotowanie danych, transformacje i augmentacje,
    \item \texttt{split\_dataset.py} --- przygotowanie struktury \texttt{train/val/test},
    \item \texttt{model.py} --- tworzenie modelu na bazie biblioteki \texttt{timm},
    \item \texttt{train.py} --- trening, checkpointy, metryki i early stopping,
    \item \texttt{predict.py} oraz \texttt{inference.py} --- inferencja dla pojedynczego obrazu,
    \item \texttt{robustness\_eval.py} --- test odporności na spadek jakości,
    \item \texttt{deepfake\_faces.py} oraz \texttt{face\_utils.py} --- wykrywanie twarzy i przygotowanie cropów,
    \item \texttt{compare\_face\_models.py} --- porównanie modelu globalnego i twarzowego,
    \item \texttt{error\_analysis.py} --- analiza błędów i eksport przykładów,
    \item \texttt{web\_demo.py} --- lokalny interfejs demonstracyjny.
\end{itemize}

Takie rozdzielenie upraszcza dalszy rozwój projektu, ponieważ etap globalny i etap twarzowy mogą być rozwijane niezależnie.

\subsection{Dane}

W eksperymentach wykorzystano zbiór danych zorganizowany w klasy \texttt{real} i \texttt{fake}, a następnie podzielony na części:

\begin{itemize}
    \item zbiór treningowy: 40800 obrazów,
    \item zbiór walidacyjny: 7200 obrazów,
    \item zbiór testowy: 12000 obrazów.
\end{itemize}

Taki podział pozwala oddzielić proces uczenia od strojenia i od finalnej oceny modelu.

W dalszym etapie przygotowano również:

\begin{itemize}
    \item curated zbiór portretów do treningu modelu twarzowego,
    \item mały benchmark OOD z nowymi generatorami GPT i Gemini / Nano2,
    \item małe zbiory adaptacyjne do prób dostrajania modeli.
\end{itemize}

\subsection{Model bazowy}

Jako model bazowy wybrano architekturę \texttt{ResNet18} \cite{he2016resnet} dostępną w bibliotece \texttt{timm}. Wybór ten wynikał z kilku powodów:

\begin{itemize}
    \item model jest lekki obliczeniowo i nadaje się do szybkich eksperymentów,
    \item posiada sprawdzoną architekturę dla klasyfikacji obrazów,
    \item można go uruchamiać efektywnie na GPU z wykorzystaniem wag wstępnie wytrenowanych.
\end{itemize}

Model przyjmuje obrazy przeskalowane do rozmiaru \texttt{224 x 224} i zwraca prawdopodobieństwo dla dwóch klas: \texttt{fake} oraz \texttt{real}.

\subsection{Etap twarzowy}

Do analizy portretów zastosowano osobny pipeline. Najpierw z obrazu wejściowego wykrywana jest twarz przy użyciu klasyfikatora Haar Cascade \cite{viola2001rapid}, a następnie wycinany jest crop w stylu \texttt{face} lub \texttt{portrait}. Tak przygotowany obraz trafia do osobnego klasyfikatora twarzowego opartego o \texttt{ConvNeXt Tiny} \cite{liu2022convnext}.

Takie rozdzielenie miało umożliwić:

\begin{itemize}
    \item odciążenie modelu od pełnego kontekstu kadru,
    \item skupienie analizy na obszarze najbardziej istotnym dla deepfake'ów,
    \item porównanie zachowania modelu globalnego i wyspecjalizowanego modelu twarzowego.
\end{itemize}

\subsection{Augmentacje i odporność}

Ze względu na charakter problemu już na etapie projektowania uznano, że zwykły trening na czystych obrazach nie będzie wystarczający. Dlatego zastosowano augmentacje związane z pogorszeniem jakości:

\begin{itemize}
    \item kompresję JPEG,
    \item rozmycie Gaussa,
    \item dodawanie szumu,
    \item operacje downscale-upscale.
\end{itemize}

Celem tych modyfikacji było uodpornienie modelu na część przypadków spotykanych w praktyce, zwłaszcza po przejściu obrazów przez komunikatory i serwisy internetowe.

\subsection{Środowisko uruchomieniowe}

Projekt jest rozwijany w Pythonie z wykorzystaniem między innymi bibliotek:

\begin{itemize}
    \item \texttt{PyTorch},
    \item \texttt{torchvision},
    \item \texttt{timm},
    \item \texttt{Pillow},
    \item \texttt{OpenCV},
    \item \texttt{Gradio}.
\end{itemize}

Trening został zoptymalizowany pod GPU poprzez:

\begin{itemize}
    \item wykorzystanie \texttt{AMP},
    \item automatyczny dobór \texttt{batch size},
    \item \texttt{channels\_last},
    \item \texttt{cudnn\_benchmark},
    \item dostrojenie \texttt{DataLoadera}.
\end{itemize}

Te decyzje projektowe pozwoliły wyraźnie poprawić praktyczne wykorzystanie karty graficznej podczas treningu.

   \newpage
\section{Implementacja}

\subsection{Trening modelu}

Najważniejszym elementem implementacji jest skrypt \texttt{train.py}, który odpowiada za:

\begin{itemize}
    \item wczytanie konfiguracji z pliku \texttt{config.yaml},
    \item przygotowanie modelu,
    \item trening i walidację,
    \item zapis najlepszych wag,
    \item wznowienie treningu z checkpointu,
    \item zapis historii uczenia i podsumowania do pliku JSON.
\end{itemize}

W trakcie pracy dodano także mechanizmy praktyczne, które okazały się istotne dla dłuższych eksperymentów:

\begin{itemize}
    \item \texttt{early stopping},
    \item checkpointy \texttt{best\_model.pt}, \texttt{last\_checkpoint.pt} i \texttt{interrupted\_checkpoint.pt},
    \item obsługę treningu na GPU z wykorzystaniem \texttt{AMP},
    \item automatyczny dobór rozmiaru batcha,
    \item lepsze ustawienia loadera pod Windows i CUDA.
\end{itemize}

\subsection{Metryki}

Podczas ewaluacji modelu raportowane są:

\begin{itemize}
    \item \texttt{Accuracy},
    \item \texttt{Precision Macro},
    \item \texttt{Recall Macro},
    \item \texttt{F1 Macro},
    \item \texttt{ROC-AUC}.
\end{itemize}

Metryki te pozwalają ocenić model nie tylko pod względem skuteczności ogólnej, ale też pod kątem jakości rozróżniania klas.

\subsection{Wyniki etapu bazowego}

Wcześniejszy model bazowy osiągnął na zbiorze testowym:

\begin{itemize}
    \item Accuracy: 0.9339,
    \item F1 Macro: 0.9339,
    \item ROC-AUC: 0.9831.
\end{itemize}

Następnie przeprowadzono drugi trening z mocniejszymi augmentacjami odpornościowymi. Wyniki tego wariantu były następujące:

\begin{itemize}
    \item Accuracy: 0.9268,
    \item Precision Macro: 0.9273,
    \item Recall Macro: 0.9267,
    \item F1 Macro: 0.9267,
    \item ROC-AUC: 0.9804.
\end{itemize}

Wariant ten okazał się minimalnie słabszy na czystym zbiorze testowym, ale lepiej przygotowany do dalszych analiz odporności.

\subsection{Etap twarzowy}

W drugim etapie projektu przygotowano osobną ścieżkę analizy portretów i cropów twarzy. W praktyce oznacza to:

\begin{itemize}
    \item automatyczne wykrycie twarzy klasyfikatorem Haar Cascade,
    \item zapis cropu w stylu \texttt{face} albo \texttt{portrait},
    \item trening osobnego modelu twarzowego niezależnego od klasyfikatora globalnego,
    \item możliwość porównania modelu globalnego i modelu twarzowego na tym samym zbiorze portretów.
\end{itemize}

Wariant twarzowy został oparty o architekturę \texttt{ConvNeXt Tiny}. Dla curated miksu wieloźródłowego uzyskano bardzo dobre wyniki na danych zbliżonych do treningowych:

\begin{itemize}
    \item Accuracy: 0.9763,
    \item Precision Macro: 0.9765,
    \item Recall Macro: 0.9763,
    \item F1 Macro: 0.9763,
    \item ROC-AUC: 0.9982.
\end{itemize}

Pokazało to, że dla klasycznych danych portretowych model twarzowy może być bardzo skuteczny, ale nie przesądza jeszcze o jego odporności na nowe generatory obrazów.

\subsection{Test odporności}

Dla aktualnego modelu przeprowadzono osobny test odporności na spadek jakości obrazu. Sprawdzono między innymi przypadki:

\begin{itemize}
    \item lekkiej i silnej kompresji JPEG,
    \item lekkiego i silnego rozmycia,
    \item lekkiego i silnego downscale,
    \item profilu mieszanego \texttt{mixed\_quality}.
\end{itemize}

Najważniejsze obserwacje były następujące:

\begin{itemize}
    \item lekkie zaburzenia obniżały skuteczność tylko nieznacznie,
    \item najsilniejszy spadek pojawiał się dla \texttt{jpeg\_extreme},
    \item przypadki mieszane obniżały skuteczność wyraźnie bardziej niż pojedyncza lekka degradacja.
\end{itemize}

Pokazuje to, że model jest już w pewnym stopniu odporny na pogorszenie jakości, ale nadal pozostaje wrażliwy na obrazy bardzo mocno skompresowane.

\subsection{Benchmark OOD dla nowych generatorów}

Aby sprawdzić, czy modele generalizują poza główny rozkład treningowy, przygotowano mały benchmark OOD złożony z:

\begin{itemize}
    \item 20 obrazów wygenerowanych przez GPT,
    \item 20 obrazów wygenerowanych przez Gemini / Nano2,
    \item 40 prawdziwych zdjęć portretowych dobranych tak, aby stylistycznie przypominały obrazy syntetyczne.
\end{itemize}

Benchmark ten nie był częścią głównego zbioru treningowego. Jego celem było możliwie uczciwe sprawdzenie, jak modele zachowują się na nowszych generatorach i bardziej współczesnych portretach. W praktyce był to więc mały test zewnętrzny, pozwalający sprawdzić nie tylko skuteczność klasyfikacji, ale także odporność na przesunięcie rozkładu danych.

Na tym benchmarku model globalny uzyskał:

\begin{itemize}
    \item Accuracy: 0.7125,
    \item Precision Macro: 0.7934,
    \item Recall Macro: 0.7125,
    \item F1 Macro: 0.6912,
    \item ROC-AUC: 0.8413.
\end{itemize}

Macierz pomyłek pokazała istotną asymetrię:

\begin{itemize}
    \item 39 z 40 obrazów \texttt{fake} zostało wykrytych poprawnie,
    \item tylko 18 z 40 obrazów \texttt{real} zostało rozpoznanych poprawnie,
    \item aż 22 prawdziwe zdjęcia zostały oznaczone jako \texttt{fake}.
\end{itemize}

Wniosek z tego etapu jest ważny: model globalny nie zawodził głównie na nowych obrazach syntetycznych, lecz na nowych realistycznych portretach, dla których generował zbyt wiele fałszywych alarmów.

\subsection{Tabela zbiorcza wyników}

Najważniejsze wyniki eksperymentów zestawiono w tabeli \ref{tab:wyniki-zbiorcze}. Pozwala ona porównać modele nie tylko na danych zbliżonych do treningowych, ale również na małym benchmarku OOD oraz po krótkiej adaptacji modelu twarzowego.

\begin{table}[H]
\centering
\caption{Zbiorcze wyniki najważniejszych eksperymentów}
\label{tab:wyniki-zbiorcze}
\begin{tabular}{|p{4.3cm}|c|c|c|}
\hline
\textbf{Wariant} & \textbf{Accuracy} & \textbf{F1 Macro} & \textbf{ROC-AUC} \\
\hline
Model globalny -- baseline test & 0.9339 & 0.9339 & 0.9831 \\
\hline
Model globalny -- wariant robust & 0.9268 & 0.9267 & 0.9804 \\
\hline
Model twarzowy \texttt{ConvNeXt Tiny} -- curated test & 0.9763 & 0.9763 & 0.9982 \\
\hline
Model globalny -- benchmark OOD & 0.7125 & 0.6912 & 0.8413 \\
\hline
Model twarzowy przed adaptacją -- OOD (wykryte twarze) & ok. 0.51 & ok. 0.49 & ok. 0.52 \\
\hline
Model twarzowy po adaptacji -- test wewnętrzny & 0.8548 & 0.8548 & 0.9469 \\
\hline
\end{tabular}
\end{table}

\subsection{Analiza jakościowa błędów}

Same metryki nie pokazują jeszcze, \emph{jakiego rodzaju} obrazy są problematyczne. Dlatego wykonano również analizę jakościową przypadków błędnych klasyfikacji.

Na rysunku \ref{fig:ood-false-positive} pokazano przykład fałszywego alarmu modelu globalnego na prawdziwym portrecie. Tego typu przypadki były częste w benchmarku OOD i sugerowały, że model nauczył się częściowo traktować estetyczne, czyste portrety jako obrazy podejrzane.

\begin{figure}[H]
\centering
\includegraphics[width=0.42\textwidth]{../reports/ood_pilot_global/test/examples/false_positive/04__gt_real__pred_fake__real_matched_real__0002__pexels-a-darmel-8158953_jpg.jpg}
\caption{Przykład false positive na benchmarku OOD: prawdziwy portret sklasyfikowany jako \texttt{fake}}
\label{fig:ood-false-positive}
\end{figure}

Na rysunku \ref{fig:ood-false-negative} pokazano odwrotny przypadek, czyli obraz syntetyczny wygenerowany przez GPT, który został oznaczony jako \texttt{real}. Takie próbki były rzadsze niż false positive modelu globalnego, ale bardzo istotne z punktu widzenia wiarygodności systemu.

\begin{figure}[H]
\centering
\includegraphics[width=0.42\textwidth]{../reports/ood_pilot_global/test/examples/false_negative/01__gt_fake__pred_real__fake_gpt__0016__gpt__p017__01_png.png}
\caption{Przykład false negative na benchmarku OOD: obraz syntetyczny sklasyfikowany jako \texttt{real}}
\label{fig:ood-false-negative}
\end{figure}

Takie przykłady potwierdzają, że problem generalizacji nie dotyczy wyłącznie pojedynczych metryk, ale ma również wymiar jakościowy: modele potrafią mylić się na wizualnie przekonujących portretach nawet wtedy, gdy na danych wewnętrznych osiągają bardzo wysokie wyniki.

\subsection{Wyniki modelu twarzowego na nowych portretach}

Porównanie wcześniejszych checkpointów twarzowych z benchmarkiem OOD pokazało, że sam dobry wynik na danych wewnętrznych nie wystarcza do poprawnej generalizacji na nowsze obrazy. Przed adaptacją:

\begin{itemize}
    \item wariant \texttt{faces\_hard\_adapt} uzyskał Accuracy około 0.48 na podzbiorze z wykrytą twarzą,
    \item wariant \texttt{faces\_convnext\_v2} uzyskał Accuracy około 0.51 na tym samym benchmarku.
\end{itemize}

Oznacza to, że pierwotna wersja modelu twarzowego nie generalizowała poprawnie na najnowsze fotorealistyczne obrazy GPT / Gemini.

W dalszym kroku wykonano mały fine-tuning modelu \texttt{faces\_convnext\_v2} na cropach twarzy przygotowanych z nowych obrazów. Na wewnętrznym teście adaptacyjnym uzyskano:

\begin{itemize}
    \item Accuracy: 0.8548,
    \item Precision Macro: 0.8552,
    \item Recall Macro: 0.8548,
    \item F1 Macro: 0.8548,
    \item ROC-AUC: 0.9469.
\end{itemize}

Wynik ten należy jednak interpretować ostrożnie, ponieważ dotyczy niewielkiego zbioru adaptacyjnego, a ręczne testy w interfejsie nadal pokazywały trudne przypadki, w których nowe obrazy syntetyczne były klasyfikowane jako \texttt{real}. Oznacza to, że adaptacja poprawiła zachowanie modelu, ale nie rozwiązała jeszcze problemu w sposób pełny.

\subsection{Interfejs demonstracyjny}

W celu zwiększenia użyteczności projektu przygotowano lokalne demo webowe w \texttt{Gradio}. Aplikacja umożliwia:

\begin{itemize}
    \item załadowanie pojedynczego obrazu,
    \item wyświetlenie predykcji \texttt{real/fake},
    \item pokazanie pewności modelu,
    \item wygenerowanie mapy Grad-CAM,
    \item pokazanie ostrzeżenia przy niskiej jakości obrazu,
    \item porównanie predykcji modelu globalnego i modelu twarzowego dla portretów.
\end{itemize}

Interfejs ten ma charakter demonstracyjny i stanowi przydatne narzędzie do ręcznej analizy zachowania modelu.

   \newpage
\section{Wnioski}

Projekt należy uznać za udany na poziomie inżynierskim i badawczym. Udało się zbudować kompletny pipeline do klasyfikacji obrazów \texttt{real} oraz \texttt{fake}, przygotować wariant globalny i twarzowy, wdrożyć narzędzia interpretowalności oraz przeprowadzić eksperymenty zarówno na wydzielonych zbiorach testowych, jak i na małym benchmarku OOD z nowymi generatorami.

Do najważniejszych osiągnięć należy zaliczyć:

\begin{itemize}
    \item poprawne uruchomienie treningu na GPU,
    \item uzyskanie wysokich wyników na bazowym zbiorze testowym,
    \item zbudowanie osobnego etapu twarzowego opartego o cropy twarzy,
    \item wdrożenie interpretowalności poprzez Grad-CAM,
    \item przygotowanie testu odporności na pogorszenie jakości,
    \item stworzenie lokalnego interfejsu demonstracyjnego,
    \item przygotowanie benchmarku OOD dla nowszych generatorów i przeprowadzenie analizy błędów.
\end{itemize}

\subsection{Co udało się osiągnąć}

Najważniejszym rezultatem projektu jest działający system, który obejmuje cały proces: przygotowanie danych, trening modeli, ewaluację, analizę odporności, interpretowalność oraz lokalny interfejs demonstracyjny. Model globalny osiągnął wysoką dokładność na głównym zbiorze testowym, a model twarzowy uzyskał bardzo dobre wyniki na danych portretowych zbliżonych do treningowych.

Udało się również pokazać, że sama dokładność na klasycznym zbiorze testowym nie jest wystarczająca do pełnej oceny systemu. Dlatego w projekcie wykonano dodatkowe testy odporności oraz benchmark OOD. Dzięki temu dokumentacja pokazuje nie tylko najlepsze wyniki, ale też ograniczenia modelu w trudniejszych warunkach.

\begin{table}[H]
\centering
\caption{Podsumowanie rezultatów i najważniejszych uwag}
\label{tab:cele-projektu}
\begin{tabular}{|L{4.2cm}|L{5.0cm}|L{4.2cm}|}
\hline
\textbf{Obszar} & \textbf{Co udało się osiągnąć} & \textbf{Ograniczenia lub uwagi} \\
\hline
Model globalny & Wysoka skuteczność na głównym zbiorze testowym: Accuracy 0.9339 dla baseline oraz 0.9268 dla wariantu robust. & Spadek skuteczności na benchmarku OOD do Accuracy 0.7125. \\
\hline
Odporność na jakość obrazu & Model zachował dobrą jakość dla lekkiego rozmycia i lekkiego downscale. & Silna kompresja JPEG oraz profil \texttt{mixed\_quality} wyraźnie obniżały wyniki. \\
\hline
Model twarzowy & Bardzo dobry wynik na wydzielonym teście portretowym: Accuracy 0.9763. & Wynik nie przenosił się automatycznie na nowe generatory spoza głównego rozkładu danych. \\
\hline
Adaptacja twarzowa & Poprawa po fine-tuningu na nowych cropach: Accuracy 0.8548 na teście adaptacyjnym. & Zbiór adaptacyjny był mały, więc wynik nie gwarantuje pełnej stabilności na dowolnych danych. \\
\hline
Interfejs demonstracyjny & Możliwość ręcznego testowania obrazów, porównania modeli i generowania Grad-CAM. & Interfejs spełnia założoną funkcję demonstracyjną i diagnostyczną. \\
\hline
\end{tabular}
\end{table}

Jednocześnie analiza praktyczna pokazała, że nawet bardzo dobry wynik ilościowy na danych zbliżonych do treningowych nie oznacza pełnej gotowości modelu do pracy na dowolnych danych. Szczególnie problematyczne okazały się:

\begin{itemize}
    \item nowe realistyczne portrety spoza głównego rozkładu treningowego,
    \item obrazy po silnej kompresji,
    \item najnowsze fotorealistyczne generatory obrazów,
    \item małe zbiory adaptacyjne, które nie pozwalają jeszcze na pełną stabilizację modelu,
    \item ograniczona dostępność dużych, publicznych i łatwo porównywalnych zbiorów obejmujących najnowsze zamknięte generatory obrazów.
\end{itemize}

\subsection{Co było problematyczne i dlaczego}

Największym problemem okazała się generalizacja, czyli przeniesienie skuteczności z wydzielonych zbiorów testowych zbliżonych do danych treningowych na nowe obrazy spoza głównego rozkładu. Przyczyną był przede wszystkim inny rozkład danych: nowe portrety i obrazy wygenerowane przez GPT / Gemini miały inną estetykę, inne artefakty oraz często wyższą jakość wizualną niż część danych treningowych.

Drugim problemem były fałszywe alarmy dla prawdziwych portretów. Model globalny bardzo dobrze wykrywał większość obrazów syntetycznych w benchmarku OOD, ale jednocześnie oznaczał wiele prawdziwych zdjęć jako \texttt{fake}. Możliwą przyczyną jest to, że model nauczył się traktować czyste, studyjne albo bardzo estetyczne portrety jako podejrzane, ponieważ podobne cechy często występują również w obrazach generowanych przez AI.

Trzecim problemem była wrażliwość na silne pogorszenie jakości obrazu. Kompresja JPEG, rozmycie i skalowanie mogą usuwać drobne artefakty, na których model opiera decyzję. W efekcie obraz po przetworzeniu przez komunikator albo serwis społecznościowy może być trudniejszy do klasyfikacji niż oryginalna próbka.

Osobnym ograniczeniem był model twarzowy. Analiza samej twarzy była przydatna, ale nie wystarczyła do pełnego rozwiązania problemu. Przyczyną jest to, że crop twarzy usuwa część kontekstu obrazu, a jednocześnie wymaga poprawnej detekcji twarzy. Dodatkowo mały zbiór adaptacyjny nie dawał wystarczająco dużo przykładów, aby model stabilnie nauczył się wszystkich wariantów nowych generatorów.

Istotnym ograniczeniem niezależnym od samej implementacji była także dostępność danych. Najnowsze modele generowania obrazów rozwijają się bardzo szybko, a publiczne zbiory referencyjne zwykle pojawiają się z opóźnieniem, są budowane dla innych celów albo nie obejmują dokładnie tych samych generatorów, stylów i typów obrazów, które są problematyczne w praktyce. Z tego powodu w projekcie nie dało się oprzeć oceny wyłącznie na dużym, gotowym i powszechnie uznanym benchmarku dla najnowszych modeli. Konieczne było przygotowanie własnego, mniejszego benchmarku OOD, który pozwolił sprawdzić kierunek problemu, ale nie zastępuje dużego zbioru referencyjnego.

Najważniejszy wniosek metodologiczny jest następujący: model globalny stanowi mocny baseline dla fotografii i znanych danych, natomiast generalizacja na najnowsze generatory GPT / Gemini pozostaje ograniczona. W benchmarku OOD problem nie polegał wyłącznie na niewykrywaniu obrazów syntetycznych, ale również na generowaniu zbyt wielu fałszywych alarmów dla nowych prawdziwych portretów.

Równie istotny jest wniosek dotyczący etapu twarzowego. Sam fakt zastosowania osobnego modelu twarzowego nie gwarantuje poprawy na nowych danych. Pierwotne checkpointy twarzowe nie generalizowały dobrze na nowe obrazy GPT / Gemini, a mała adaptacja na nowych cropach poprawiła wynik na teście adaptacyjnym. Nadal jednak ręczne testy pokazały trudne przypadki, w których nawet adaptowany model twarzowy klasyfikował syntetyczny portret jako \texttt{real}.

Oznacza to, że końcowy efekt pracy należy interpretować jako:

\begin{itemize}
    \item działający system badawczy z pełnym pipeline'em,
    \item mocny baseline dla starszych i znanych danych,
    \item narzędzie pozwalające badać ograniczenia detekcji na nowszych generatorach,
    \item punkt wyjścia do dalszych prac nad lepszą generalizacją modeli twarzowych.
\end{itemize}

W kolejnych pracach najbardziej uzasadnione byłyby:

\begin{enumerate}
    \item rozbudowanie zbioru OOD o większą liczbę obrazów z najnowszych generatorów, najlepiej po pojawieniu się większych publicznych benchmarków dla tych modeli,
    \item przygotowanie większego i bardziej zrównoważonego zbioru adaptacyjnego dla twarzy,
    \item sprawdzenie silniejszych metod detekcji twarzy i alternatywnych strategii cropowania,
    \item dalsza analiza modeli hybrydowych, które łączą sygnał globalny i twarzowy.
\end{enumerate}

Można zatem stwierdzić, że projekt osiągnął wyraźny i mierzalny progres: posiada sprawdzony etap bazowy, rozszerzenie twarzowe, benchmark OOD oraz udokumentowaną analizę ograniczeń. Z punktu widzenia pracy końcowej jest to wartościowy rezultat, ponieważ pokazuje nie tylko implementację, ale też realistyczną ocenę granic skuteczności obecnych metod detekcji.


%%%%%%%%%%%%%%%%%%% koniec treści głównej dokumentu %%%%%%%%%%%%%%%%%%%%%

    \newpage
    \addcontentsline{toc}{section}{Literatura}
    \printbibliography

    \newpage
    \hypersetup{linkcolor=black}
    \renewcommand{\cftparskip}{3pt}
    \clearpage
    \renewcommand{\cftloftitlefont}{\Large\bfseries\sffamily}
    \listoffigures
    \addcontentsline{toc}{section}{Spis rysunków}
    \thispagestyle{fancy}

    \newpage
    \renewcommand{\cftlottitlefont}{\Large\bfseries\sffamily}
    \def\listtablename{Spis tabel}
    \addcontentsline{toc}{section}{Spis tabel}\listoftables
    \thispagestyle{fancy}

    \newpage
    \renewcommand{\cftlottitlefont}{\Large\bfseries\sffamily}
    \renewcommand\lstlistlistingname{Spis listingów}
    \addcontentsline{toc}{section}{Spis listingów}\lstlistoflistings
    \thispagestyle{fancy}

    % lista rzeczy do zrobienia: wypisuje na końcu dokumentu, patrz pakiet todo.sty
    \todos
\end{document}
